% answers_short.tex —— 1977 年·江西卷 · 简短答案速查
% 由 tools/make_short_answers.py 自 answers.tex 的 \gksolhead/\gkanswer 提取，
% 并与 1977/1977_数学真题_江西_答案要点.md 逐题比对后生成，勿手改（改源 MD 后重跑）。

\begin{gknotice}
\textbf{使用说明}：本节只给最终答案便于快速核对，完整推导见解析版。答案为编写者独立推导，\textbf{非官方标准答案}。
\end{gknotice}

\gkpart{题目}
\begin{itemize}[leftmargin=2.6em]
  \item[\textbf{1.}] (1) $-5$；(2) $\dfrac{\pi}{4}$；(3) $2$；(4) $4$；(5) $-1$；(6) 模 $1$，幅角 $\dfrac{\pi}{3}$，三角函数式 $\cos\dfrac{\pi}{3}+i\sin\dfrac{\pi}{3}$；(7) $\dfrac{\sqrt2}{2}$；(8) 小（$ac<bc$）；(9) 中；(10) $8$；(11) $-3$；(12) $4$ 公分；(13) 振幅 $2$，周期 $\dfrac{2\pi}{3}$；(14) 等差中项 $5$，等比中项 $\pm1$；(15) 正方形边长 $\sqrt2\,R$，正三角形边长 $\sqrt3\,R$
  \item[\textbf{2.}] 最大面积为 $\dfrac{36}{5}$（面积单位），此时两邻边均为 $3$。
  \item[\textbf{3.}] $-\dfrac{243}{8}=-30\dfrac{3}{8}$。
  \item[\textbf{4.}] (1) $72$ 个；(2) $18$ 个
  \item[\textbf{5.}] 见解析版
  \item[\textbf{6.}] $V=\dfrac{128}{9}(2-\sqrt3)$ 立方公分。
\end{itemize}
